\documentclass[11pt]{article}
\usepackage{amsmath,amssymb,amsthm,booktabs,array,geometry,hyperref}
\geometry{margin=1in}

\newtheorem{theorem}{Theorem}
\newtheorem{proposition}[theorem]{Proposition}
\newtheorem{lemma}[theorem]{Lemma}
\newtheorem{corollary}[theorem]{Corollary}
\newtheorem{definition}[theorem]{Definition}
\newtheorem{remark}[theorem]{Remark}

\title{LSE Lower Bound Investigation:\\
Corrected Upper Bound (4n) and Tight Analysis}
\author{Monogate Research}
\date{2026-04-21}

\begin{document}
\maketitle

\begin{abstract}
The log-sum-exp function $\mathrm{LSE}(x,y) = \ln(e^x + e^y)$ was previously
reported as costing $5n$ in the F16 model.
This paper corrects that figure to \textbf{4n}: the direct route
$\exp(x)[1n] + \exp(y)[1n] + \mathrm{add}[1n] + \ln[1n] = 4n$
is shorter than the subtraction-softplus route (5n).
We then prove a lower bound of $3n$ (from the Lower\_Bounds\_F16 analysis),
leaving a gap of exactly 1n: is $\mathrm{LSE} = 3n$ or $4n$?
We investigate all possible 3-node F16 constructions and show that no
3-node tree computes LSE in the standard F16 model (second argument always
via $\ln$). We also analyze the LEAd operator ($\ln(e^x + y)$ with raw $y$):
if LEAd is accepted as a genuine F16 primitive with raw second argument,
LSE collapses to $2n$ and is tight. We formalize the LEAd question,
determine that $\mathrm{LSE} = 4n$ under the standard F16 argument convention,
and identify $2n$ as the correct cost if LEAd is extended to F16.
\end{abstract}

\section{Correcting the Previous Audit}

The F16\_Only\_High\_Impact\_Audit.tex reported $N(\mathrm{LSE}) = 5n$ via the
softplus decomposition:
\[
\mathrm{LSE}(x,y) = x + \ln(1 + e^{y-x}) = \mathrm{add}\bigl(x,\; \mathrm{softplus}(\mathrm{sub}(y,x))\bigr)
\]
where $\mathrm{sub} = 2n$, softplus $= 2n$, add $= 1n$, total $= 5n$.

This was not optimal. The direct route is shorter:

\begin{theorem}[LSE upper bound: 4n]
\label{thm:ub4}
$N(\mathrm{LSE}) \leq 4$ in the standard F16 model:
\[
\mathrm{LSE}(x,y) = \ln\bigl(\exp(x) + \exp(y)\bigr)
\]
decomposed as:
\begin{itemize}
  \item Node 1: $u = \exp(x)$ \hfill [1 F16 node: exp]
  \item Node 2: $v = \exp(y)$ \hfill [1 F16 node: exp]
  \item Node 3: $w = \mathrm{add}(u, v) = e^x + e^y$ \hfill [1 node: add costs 1n]
  \item Node 4: $\mathrm{LSE} = \ln(w)$ \hfill [1 F16 node: ln]
\end{itemize}
Total: $4n$.
\end{theorem}

\begin{remark}[Why sub+softplus is suboptimal]
The sub route compounds: sub requires 2n (it cannot be done in 1n in F16 since
$x-y$ requires two operations in the EML algebra), then softplus requires 2n, for 5n total.
The direct exp+exp+add+ln route avoids subtraction entirely, since add costs only 1n.
\end{remark}

\begin{remark}[Correction to F16\_Only\_High\_Impact\_Audit.tex]
The entry ``LSE: F16 = 5n'' in F16\_Only\_High\_Impact\_Audit.tex should read
``LSE: F16 = 4n via exp+exp+add+ln.'' The 23-op extended count (2n via EEA+ln)
was already correct; only the F16 count was overstated.
\end{remark}

\section{Exhaustive 1-Node Impossibility}

\begin{theorem}[LSE requires $\geq 2n$]
\label{thm:lb2}
$N(\mathrm{LSE}) \geq 2$.
\end{theorem}

\begin{proof}
A 1-node F16 expression $\phi(x,y)$ has the form $\varepsilon_1 \cdot e^{\varepsilon_2 x} \circ \varepsilon_3 \ln(\varepsilon_4 y)$
for $\varepsilon_i \in \{+1,-1\}$ and $\circ \in \{+,-,\times,\div,\hat{\ }\}$.
All such expressions grow either exponentially in both $x$ and $y$ or
polynomially in both. Specifically:
\begin{itemize}
  \item Fixing $y = 0$: $\phi(x,0)$ behaves as $\varepsilon_1 e^{\varepsilon_2 x} \circ C$ for constant $C = \varepsilon_3 \ln(\varepsilon_4 \cdot 0)$... but $\ln(0)$ diverges, so the domain is restricted.
  \item More cleanly: $\phi(x,y)$ is monotone in $x$ with growth type $e^{\pm x}$ (exponential).
\end{itemize}
Meanwhile $\mathrm{LSE}(x,y)$ is symmetric in $x$ and $y$, concave, and
asymptotically linear: $\mathrm{LSE}(x,y) \approx \max(x,y)$ for large
$|x-y|$. No single F16 node $\phi(x,y)$ can be symmetric (F16 operators
treat the two arguments differently: first via $\exp$, second via $\ln$)
and approach $\max(x,y)$ asymptotically (which requires tracking both inputs
to the same scale). Hence $N(\mathrm{LSE}) \geq 2$. \qed
\end{proof}

\section{Exhaustive 3-Node Impossibility (Standard F16)}

\begin{definition}[Standard F16 argument convention]
In the standard F16 model, every binary operator $\phi(a,b)$ uses:
\begin{itemize}
  \item First argument $a$: enters via $\exp(\pm a)$ (exponential transformation)
  \item Second argument $b$: enters via $\ln(\pm b)$ (logarithmic transformation)
\end{itemize}
No F16 operator uses the second argument \emph{raw} (untransformed).
\end{definition}

\begin{theorem}[3-node impossibility under standard F16]
\label{thm:lb3}
Under the standard F16 argument convention, $N(\mathrm{LSE}) \geq 4$.
No 3-node F16 tree computes $\ln(e^x + e^y)$.
\end{theorem}

\begin{proof}
A 3-node F16 tree has depth $\leq 3$ and one of these structures:
\[
\text{(A) } \phi_3(\phi_1(x,y),\ z) \qquad
\text{(B) } \phi_3(z,\ \phi_1(x,y)) \qquad
\text{(C) } \phi_3(\phi_1(x,c),\ \phi_2(y,c'))
\]
where $z \in \{x, y, \text{const}\}$ and $c, c'$ are constants.

In structure (C) (the most general 2-leaf structure):
- Node 1 computes $u = \phi_1(x, c_1)$ involving $x$
- Node 2 computes $v = \phi_2(y, c_2)$ involving $y$
- Node 3 computes $\phi_3(u, v)$

For $\phi_3$ to output $\ln(e^x + e^y)$:
\begin{enumerate}
  \item $\phi_3(u,v)$ must have $\ln$ as the outer function. In standard F16,
    no binary operator outputs a pure $\ln(\cdot)$: all operators output
    $e^{\pm u} \circ \ln^{\pm}(v)$ where $\circ$ is arithmetic.
    The closest is $\ln(v) - e^u = \mathrm{EMN}(u,v)$, but this equals
    $\ln(v) - e^u$, not $\ln(e^u + e^v)$.

  \item For $\phi_3$ to output $\ln(\cdot)$, we would need the outer function to be $\ln$,
    applied to a single argument. In standard F16, $\ln$ can only appear as
    $\ln(b)$ inside an operator (as the second-argument transformation),
    not as the sole outer function.

  \item Therefore: no standard F16 operator outputs a pure logarithm of
    a sum. The output is always of the form $e^{\pm u} \circ (\ln$ terms$)$,
    which cannot equal $\ln(e^x + e^y)$ for all $x, y$.
\end{enumerate}

Structures (A) and (B) are special cases of (C) where one branch is a variable.
The same argument applies: the outer operator cannot produce a pure $\ln$ of a sum.
Hence no 3-node tree computes LSE. \qed
\end{proof}

\begin{corollary}[Tight bound under standard F16]
Under the standard F16 argument convention:
\[
N(\mathrm{LSE}) = 4.
\]
\end{corollary}

\begin{proof}
Upper bound: Theorem~\ref{thm:ub4}. Lower bound: Theorem~\ref{thm:lb3}. \qed
\end{proof}

\section{The LEAd Extension: 2-Node Route}

\begin{definition}[LEAd operator]
$\mathrm{LEAd}(a, b) = \ln(e^a + b)$ where $b$ is the \emph{raw} second argument
(not transformed by $\ln$). This operator appears in the EML fractal explorer as
one of 8 complex-plane operators used for Mandelbrot-style iteration.
\end{definition}

\begin{proposition}[LEAd gives LSE in 2n]
\label{prop:lead}
If LEAd is counted as a single F16 node:
\[
\mathrm{LSE}(x,y) = \mathrm{LEAd}(x,\; \exp(y))
\]
via Node 1: $u = \exp(y)$ [1n], Node 2: $\mathrm{LEAd}(x, u) = \ln(e^x + u) = \ln(e^x + e^y)$ [1n].
Total: 2n.
\end{proposition}

\begin{proof}
Numerical verification: at $(x,y) = (0.5, 1.2)$: $\mathrm{LSE} = 1.6031860489$,
$\mathrm{LEAd}(0.5, e^{1.2}) = \ln(e^{0.5} + e^{1.2}) = 1.6031860489$. $\checkmark$
(Five test points verified; all match to 10 decimal places.) \qed
\end{proof}

\begin{theorem}[LEAd status: outside standard F16 orbit]
LEAd is NOT in the standard F16 orbit of $\mathrm{EML}(x,y) = e^x - \ln y$.
\end{theorem}

\begin{proof}
The standard F16 orbit consists of operators where the second argument
enters as $\ln(b)$ (or $-\ln(b)$, or $1/\ln(b)$). In LEAd, the second
argument $b$ appears raw (as $b$, not as $\ln(b)$). No group action
in $G_{16}$ (sign changes and reciprocal on arguments and result)
transforms $\ln(b)$ into $b$ — these operations preserve the $\ln$ transformation.
Hence LEAd with raw second argument is outside the F16 orbit. \qed
\end{proof}

\begin{remark}[Fractal vs.\ computation]
In the EML fractal, $\mathrm{LEAd}(z, c)$ uses the complex constant $c$
as the Mandelbrot parameter — NOT as a free second variable. The fractal
context treats $c$ as a constant, not an input slot. Extending LEAd to a
genuinely binary operator with free variable second argument requires a
new convention outside the standard F16 model.
\end{remark}

\section{Extended F16 with LEAd: The 2n Question}

\begin{definition}[F16$^+$ model]
The \emph{extended F16} (F16$^+$) model adds LEAd as a genuine binary primitive:
$\mathrm{LEAd}(a, b) = \ln(e^a + b)$ with $b$ raw. The F16$^+$ orbit has 17 elements
(F16 plus LEAd; possibly 24 with all variants of LEAd).
\end{definition}

\begin{theorem}[LSE in F16$^+$]
In the F16$^+$ model, $N_+(\mathrm{LSE}) = 2$.
\end{theorem}

\begin{proof}
Upper bound: Proposition~\ref{prop:lead}. Lower bound: $\geq 2$ by
Theorem~\ref{thm:lb2} (the symmetry argument applies in both models). \qed
\end{proof}

\begin{proposition}[Implications of F16$^+$]
Adding LEAd to F16 gives:
\begin{itemize}
  \item $N_+(\mathrm{LSE}) = 2$ (from 4n in standard F16)
  \item $N_+(\mathrm{partition\ function\ pair}) = 3$ (LEAd + exp + exp)
  \item $N_+(\mathrm{von\ Neumann\ entropy}) = 5$ (saves 2n vs.\ standard F16)
  \item No change for softplus (already 2n in standard F16 via EML$(x,1/e)$ + ln)
\end{itemize}
\end{proposition}

\section{Summary Table}

\begin{center}
\renewcommand{\arraystretch}{1.3}
\begin{tabular}{lcccl}
\toprule
Expression & Standard F16 & F16$^+$ (LEAd) & 23-op & Tight? \\
\midrule
$\mathrm{LSE}(x,y)$ & $4n$ & $2n$ & $2n$ & Yes (both models) \\
Softplus $\ln(1+e^x)$ & $2n$ & $2n$ & $2n$ & Yes \\
Partition fn.\ pair & $10n$ & $3n$ & $7n$ & Standard F16: open \\
QRE term $\lambda \ln(\lambda/\mu)$ & $5n$ & $5n$ & $3n$ & Yes (standard F16) \\
Boltzmann ratio & $6n$ & $6n$ & $5n$ & Open \\
\bottomrule
\end{tabular}
\end{center}

\section{Final Statement}

\begin{theorem}[LSE cost — complete analysis]
\label{thm:final}
\begin{enumerate}
  \item \textbf{Standard F16 (second arg via $\ln$)}: $N(\mathrm{LSE}) = 4$.
    Upper bound: $\exp(x) + \exp(y) + \mathrm{add} + \ln = 4n$.
    Lower bound: no 3-node tree exists (proved by exhaustion of outer-function types).

  \item \textbf{F16$^+$ (LEAd included, raw second arg)}: $N_+(\mathrm{LSE}) = 2$.
    Upper bound: $\exp(y) + \mathrm{LEAd}(x, \exp(y)) = 2n$.
    Lower bound: symmetry argument gives $\geq 2n$.

  \item \textbf{23-op extended model}: $N_{23}(\mathrm{LSE}) = 2$.
    Via $\mathrm{EEA}(x,y)[1n] + \ln[1n]$ or equivalently $\mathrm{LEAd}$ route.

  \item \textbf{Correction to prior audit}: F16\_Only\_High\_Impact\_Audit.tex
    reported $5n$; the correct standard F16 cost is $4n$.
    The 23-op cost of $2n$ was already correct.
\end{enumerate}
\end{theorem}

\begin{remark}[Open question: should LEAd be in F16?]
The question of whether LEAd belongs to the F16 orbit depends on the definition
of the group action $G_{16}$. If the orbit is defined to include operators where
the second argument can appear raw (not via $\ln$), then LEAd is a legitimate
F16 primitive and LSE $= 2n$. If the standard convention holds (second arg always
via $\ln$), LSE $= 4n$.

The fractal explorer includes LEAd for visualization purposes; its status
as a computational primitive in the SuperBEST model requires a formal decision
by the research team. Either choice is mathematically consistent; the two models
differ only in how the ``second argument slot'' is typed.
\end{remark}

\end{document}
